%%
%% Bismarck Gomes <bismarck@puc-rio.br>
%%

\documentclass[mestrado]{TemplatePUC}


%%%
%%% Pacotes
%%%
\usepackage{multirow} % mesclar linhas na tabela
\usepackage{multicol} % mesclar colunas na tabela
\usepackage{lipsum}   % gerador de lero-lero
\usepackage{csquotes} % evita mensagem de erro no babel
\usepackage[style=abnt, justify]{biblatex}


%%%
%%% Referencias biliograficas
%%%

\addbibresource{bibliografia.bib}


%%%
%%% Definições globais
%%%

% Exibição de listas (figuras, tabelas e nomenclatura)
% Ex:\imprimirlista{fig}            % lista de figs e tabs 
% Ex:\imprimirlista{fig}{nom}       % lista de figs e nom 
% Ex:\imprimirlista{fig}{tab}{nom}  % lista de figs, tabs e nom 
\imprimirlista{fig}{tab}{nom}

% Numeração de figuras e tabelas por capítulo
\listasporcapitulo


%%%
%%% Titulos
%%%

\titulo{Acoplamento Termo-Hidro-Mecânico Iterativo Usando Malhas Não Correlacionadas}
%% \subtitulo{Aqui vai o subtitulo caso precise}

\autor{Bismarck Gomes Souza Júnior}
\autorR{Souza Júnior, Bismarck Gomes}

\orientador{Sergio Augusto Barreto da Fontoura}{PhD.}
\orientadorR{Fontoura, Sergio A. B.}

\coorientador{Nelson Inoue}{Dr.}
\coorientadorR{Inoue, Nelson}
\coorientadorInst{Grupo de Tecnologia e Engenharia de Petróleo}{PUC-Rio}

\dia{18}
\mes{dezembro}
\ano{2017}

% \cidade{Rio de Janeiro}
\departamento{Engenharia Civil e Ambiental}
\programa{Engenharia Civil}
\centro{Centro Técnico Científico}


%%%
%%% Banca
%%%

\banca{%
  \membrodabanca{Fulano de Tal}{PhD.}
    {Departamento de Engenharia Civil e Ambiental}{PUC-Rio}

  \membrodabanca{Beltrano Matos Figueira}{PhD.}
    {Departamento de Engenharia Civil e Ambiental}{PUC-Rio}

  \membrodabanca{Ciclano Almeida Pinto}{Dr.}
    {Departamento de Engenharia Civil}{UFRJ}

  \coordenador{José Eugenio Leal}{Prof.}
}


%%%
%%% Currículo
%%%

\curriculo{%
  Graduou-se em Engenharia de Petróleo pela Universidade Estadual do Norte Fluminense (UENF), tendo desenvolvido trabalho científico na área de simulação numérica de reservatórios. Cursou mestrado em Engenharia de Reservatório e Exploração da UENF, com uma dissertação na área de simulação numérica multifásica e composicional de poços de petróleo. Durante o período de doutorado, trabalhou junto ao Grupo de Tecnologia e Engenharia de Petróleo da PUC-Rio, onde atualmente se enquadra como pesquisador na área de simulação numérica aplicada à geomecânica de reservatórios.}


%%%
%%% Ficha catalográfica.
%%%

%% Documento colorido?
\colorido{true}  

%% Classificação Decimal de Dewey (ver com a biblioteca)
\CDD{624}

%% Prechaves
\prechaves{
  \prechave{Engenharia civil}
}

%% Chaves para o resumo e ficha catalográfica
\chaves{
  \chave{geomecânica de reservatórios}
  \chave{acoplamento termo-hidro-mecânico}
  \chave{acoplamento de malhas}
}

%% Chaves para o abstract
\chavesus{
  \chave{reservoir geomechanics}
  \chave{thermo-hydro-mechanical coupling}
  \chave{mesh coupling}
}


%%%
%%% Agradecimentos
%%%

\agradecimentos{
  \lipsum[1-3]
}


%%%
%%% Resumo
%%%
\resumo{%
  \lipsum[4]
 }


%%%
%%% Abstract
%%%

\titulous{Iterative Thermal-Hydro-Mechanical Coupling Using Uncorrelated Meshes}

\resumous{ 
  \lipsum[5]
}


%%%
%%% Epigrafe
%%%

\epigrafe{%
  \lipsum[6]
}
\epigrafeautor{Winston Churchill}
% \epigrafelivro{Regards sur le passé}


%%%
%%% Documento
%%%

\begin{document}

  %!TEX root = ./Masters.tex

%%
%% Bismarck Gomes <bismarck@puc-rio.br>
%%

\chapter{Introdução} 
\label{chap:Introducao}

%%%%%%%%%%%%%%%%%%%%%%%%%%%%%%%%%%%%%%%%%%%%%%%%%%%%%%%%%%%%%%%%%%%%%%%%%%%%%%%%
\section{Escopo do Problema} 
\label{sec:Escopo-do-Problema}

\lipsum[7-8]


%%%%%%%%%%%%%%%%%%%%%%%%%%%%%%%%%%%%%%%%%%%%%%%%%%%%%%%%%%%%%%%%%%%%%%%%%%%%%%%%
\section{Objetivos}
\label{sec:Objetivos}

\lipsum[10-12]


%%%%%%%%%%%%%%%%%%%%%%%%%%%%%%%%%%%%%%%%%%%%%%%%%%%%%%%%%%%%%%%%%%%%%%%%%%%%%%%%
\section{Organização do Documento}
\label{sec:Organização-do-Documento}

\lipsum[13-15]


  %!TEX root = ./PreThesis.tex

%%
%% Bismarck Gomes <bismarck@puc-rio.br>
%%

\chapter{Revisão Bibliográfica}
\label{chap:Revisão-Bibliográfica}

\lipsum[15]

%%%%%%%%%%%%%%%%%%%%%%%%%%%%%%%%%%%%%%%%%%%%%%%%%%%%%%%%%%%%%%%%%%%%%%%%%%%%%%%%
\section{Artigos Tecnológicos}
\label{sec:Artigos-Tecnológicos}

\lipsum[16-20]


%%%%%%%%%%%%%%%%%%%%%%%%%%%%%%%%%%%%%%%%%%%%%%%%%%%%%%%%%%%%%%%%%%%%%%%%%%%%%%%%
\section{Artigos Científicos}
\label{sec:Artigos-Científicos}

Donec molestie, magna ut luctus ultrices, tellus arcu nonummy velit, sit
amet pulvinar elit justo et mauris. In pede. Maecenas euismod elit eu erat.
Aliquam augue wisi, facilisis congue, suscipit in, adipiscing et, ante. In justo.
Cras lobortis neque ac ipsum. Nunc fermentum massa at ante. Donec orci
tortor, egestas sit amet, ultrices eget, venenatis eget, mi. Maecenas vehicula
leo semper est. Mauris vel metus. Aliquam erat volutpat. In rhoncus sapien ac
tellus. Pellentesque ligula \cite{biot1957}.

Cras dapibus, augue quis scelerisque ultricies, felis dolor placerat sem, id
porta velit odio eu elit. Aenean interdum nibh sed wisi. Praesent sollicitudin
vulputate dui. Praesent iaculis viverra augue. Quisque in libero. Aenean
gravida lorem vitae sem ullamcorper cursus. Nunc adipiscing rutrum ante.
Nunc ipsum massa, faucibus sit amet, viverra vel, elementum semper, orci.
Cras eros sem, vulputate et, tincidunt id, ultrices eget, magna. Nulla varius
ornare odio. Donec accumsan mauris sit amet augue. Sed ligula lacus, laoreet
non, aliquam sit amet, iaculis tempor, lorem. Suspendisse eros. Nam porta,
leo sed congue tempor, felis est ultrices eros, id mattis velit felis non metus.
Curabitur vitae elit non mauris varius pretium. Aenean lacus sem, tincidunt
ut, consequat quis, porta vitae, turpis. Nullam laoreet fermentum urna. Proin
iaculis lectus \cite{terzaghi1943}.

Sed mattis, erat sit amet gravida malesuada, elit augue egestas diam,
tempus scelerisque nunc nisl vitae libero. Sed consequat feugiat massa. Nunc
porta, eros in eleifend varius, erat leo rutrum dui, non convallis lectus orci
ut nibh. Sed lorem massa, nonummy quis, egestas id, condimentum at, nisl.
Maecenas at nibh. Aliquam et augue at nunc pellentesque ullamcorper. Duis
nisl nibh, laoreet suscipit, convallis ut, rutrum id, enim. Phasellus odio. Nulla
nulla elit, molestie non, scelerisque at, vestibulum eu, nulla. Ut odio nisl,
facilisis id, mollis et, scelerisque nec, enim. Aenean sem leo, pellentesque sit
amet, scelerisque sit amet, vehicula pellentesque, sapien \cite{anderson1989}.

\lipsum[40]


%%%%%%%%%%%%%%%%%%%%%%%%%%%%%%%%%%%%%%%%%%%%%%%%%%%%%%%%%%%%%%%%%%%%%%%%%%%%%%%%
\section{Crítica aos Trabalho Existentes}
\label{sec:Crítica-aos-Trabalho-Existentes}

\lipsum[25-29]

  %!TEX root = ./Thesis.tex

%%
%% Bismarck Gomes <bismarck@puc-rio.br>
%%

\chapter{Metodologia}
\label{chap:Metodologia}

Conforme discutido no Capítulo \ref{chap:Revisão-Bibliográfica} \ldots
\lipsum[1]
%%%%%%%%%%%%%%%%%%%%%%%%%%%%%%%%%%%%%%%%%%%%%%%%%%%%%%%%%%%%%%%%%%%%%%%%%%%%%%%%
\section{Figuras}
\label{sec:Figuras}

A Figura \ref{fig:puc.png} ilustra um exemplo de como referenciar uma figura, 
enquanto que a Figura \ref{fig:puc2}, como utilizar subfiguras.

\begin{figure}[ht]
  \centering
  \includegraphics[scale=1.0]{figs/puc.png}
  \caption{Logo da PUC-Rio}
  \label{fig:puc.png}
 \end{figure}

\begin{figure}[ht]
  \centering
  \subfigure[Legenda 1]{\includegraphics[scale=0.3]{figs/puc.png}}
  \hspace{1cm}
  \subfigure[Legenda 2]{\includegraphics[scale=0.4]{figs/puc.png}}
  \caption{Logo da PUC-Rio em forma de subfigura}
  \label{fig:puc2}
 \end{figure}

%%%%%%%%%%%%%%%%%%%%%%%%%%%%%%%%%%%%%%%%%%%%%%%%%%%%%%%%%%%%%%%%%%%%%%%%%%%%%%%%
\section{Equações}
\label{sec:Equações}
 
Exemplos de equações:

\begin{itemize}
	\item{Equação simples: eq. (\ref{eq:simples})}
	\item{Equação grande: eq. (\ref{eq:grande})}
	\item{Equação alinhada: eq. (\ref{eq:alinhada})}
\end{itemize}

Seguem as equações, começando com uma relação bem simples entre os lados de um triângulo retângulo escaleno:
\begin{equation} \label{eq:simples}
  a^2 = b^2 + c^2,
 \end{equation}
onde $a$ é a hipotenusa; $b$, o menor cateto; e $c$, o maior cateto.
\nomenclature[A]{$a$}{Hipotenusa}
\nomenclature[A]{$b$}{Menor cateto}
\nomenclature[A]{$c$}{Maior cateto}

Já uma equação maior, que ocupa mais de uma linha, é necessário um tratamento especial. Por exemplo, a série de Taylor de uma função $f$ em torno do ponto $x=x_0$ dada por:
\begin{multline} \label{eq:grande}
  f(x) = f(x_0)+f'(x_0)\frac{\left(x-x_0\right)}{1!} +
  			 f''(x_0)\frac{\left(x-x_0\right)^2}{2!} + \\
  			 f'''(x_0)\frac{\left(x-x_0\right)^3}{3!} 
  			 +\ldots+f^{(n)}(x_0)\frac{\left(x-x_0\right)^n}{n!}.
	\nomenclature[A]{$f$}{Função qualquer}
 \end{multline}

Por outro lado, o desenvolvimento de uma equação pode ser escrita na forma:
\begin{align} \label{eq:alinhada}
	A &= \pi r^2     \nonumber\\
	  &= \pi (10)^2  \nonumber\\
	  &= 100\pi 
  \end{align}

Da mesma forma, é possível escrever um sistema de equações:
\begin{eqnarray} \label{eq:tmp74}
    x+y+z &=& 1 \\
    3x-2y-4z &=& 10 \\
    3x-2y &=& 20
 \end{eqnarray}
ou ainda:

 \begin{equation} \label{eq:tmp81}
  \begin{cases}
    x+y +z= 1 \\
    3x-2y -4z= 10 \\
    3x-2y = 20
  \end{cases}
 \end{equation}


%%%%%%%%%%%%%%%%%%%%%%%%%%%%%%%%%%%%%%%%%%%%%%%%%%%%%%%%%%%%%%%%%%%%%%%%%%%%%%%%
\section{Tabelas}
\label{sec:Tabelas}

Seguem o exemplo de duas tabelas, Tabela \ref{tab:Primeira-tabela} e Tabela \ref{tab:Segunda-tabela}.

\begin{table}[ht]
	\centering
	\caption{Primeira tabela}
	\label{tab:Primeira-tabela}
	% \vspace{0.5em}
	\begin{tabular}{ clc } 
	\hline
	Células 	& Lugar    & Valor \\ 
	\hline
	c1   			& Caxias   & 90,0 \\ 
	c2 				& Mesquita & 59,0 \\ 
	\hline
	\end{tabular}
\end{table}

\begin{table}[ht]
	\centering
	\caption{Segunda tabela}
	\label{tab:Segunda-tabela}
	% \vspace{0.5em}
	\begin{tabular}{ clc } 
	\hline
	Células 	& Lugar    & Valor \\ 
	\hline
	\multirow{2}{1cm}{\centering RJ} 
	    			& Caxias   & 90,0 \\ 
	   				& Mesquita & 59,0 \\ 
	\hline
	\end{tabular}
\end{table}

\begin{table}[ht]
	\centering
	\caption{Segunda tabela}
	% \vspace{0.5em}
	\begin{tabular}{ clc } 
	\hline
	Células 	& Lugar    & Valor \\ 
	\hline
	\multirow{2}{1cm}{\centering RJ} 
	    			& Caxias   & 90,0 \\ 
	   				& Mesquita & 59,0 \\ 
	\hline
	\end{tabular}
\end{table}

\begin{table}[ht]
	\centering
	\caption{Segunda tabela}
	% \vspace{0.5em}
	\begin{tabular}{ clc } 
	\hline
	Células 	& Lugar    & Valor \\ 
	\hline
	\multirow{2}{1cm}{\centering RJ} 
	    			& Caxias   & 90,0 \\ 
	   				& Mesquita & 59,0 \\ 
	\hline
	\end{tabular}
\end{table}

\begin{table}[ht]
	\centering
	\caption{Segunda tabela}
	% \vspace{0.5em}
	\begin{tabular}{ clc } 
	\hline
	Células 	& Lugar    & Valor \\ 
	\hline
	\multirow{2}{1cm}{\centering RJ} 
	    			& Caxias   & 90,0 \\ 
	   				& Mesquita & 59,0 \\ 
	\hline
	\end{tabular}
\end{table}


%%%%%%%%%%%%%%%%%%%%%%%%%%%%%%%%%%%%%%%%%%%%%%%%%%%%%%%%%%%%%%%%%%%%%%%%%%%%%%%%
\section{Referências}
\label{chap:Referências}

Segundo \textcite{anderson1989}, a terra é levemente arredondada. E de acordo com \textapud{biot1957}{righetto2016} é possível chegar ao Japão, cavando um buraco bem profundo no Brazil \cite{aavatsmark2007}.

Porém, lembre-se de sempre vericar a fonte\footcite{edwards1998} de suas informações para não cometer o equívoco\footnote{Texto descritivo da seleção} de divulgar uma informação errada \apud{terzaghi1943}{zienkiewicz1999}.
\iffalse
\fi
  
  \printbibliography

  \appendix
  %!TEX root = ./Masters.tex

%%
%% Bismarck Gomes <bismarck@puc-rio.br>
%%

\chapter{Arquivos de Entrada}
\label{chap:Arquivos-de-Entrada}

\lipsum[30]


%%%%%%%%%%%%%%%%%%%%%%%%%%%%%%%%%%%%%%%%%%%%%%%%%%%%%%%%%%%%%%%%%%%%%%%%%%%%%%%%
\section{Simulação A}
\label{sec:Simulação-A}

\lipsum[31-33]


%%%%%%%%%%%%%%%%%%%%%%%%%%%%%%%%%%%%%%%%%%%%%%%%%%%%%%%%%%%%%%%%%%%%%%%%%%%%%%%%
\section{Simulação B}
\label{sec:Simulação-B}

\lipsum[35-40]


\end{document}
